% Onco-specific Intro paragraph (~140w). For Tumor Biology agent-mozurh1r.
% Complements tumor_bio_intro_para.tex (three-property framing: low abundance + adjacency + non-exchangeable cancer batches).
% Philosophy-level: WHY cancer is the hardest exchangeability case, connecting to Lemma S1.
% Target: sections/01_introduction.tex, to follow tumor_bio_intro_para.tex immediately.

Cancer is the limiting case for the exchangeability bound \Cref{lem:exchangeability}. In non-malignant tissue atlases, batches typically vary along technical axes (donor, chemistry, centre) that are approximately independent of the biology under study, so the bound is tight. Cancer atlases invert that relationship: each batch encodes a disease state --- tumour subtype, treatment history, micro-environmental niche --- as an inseparable covariate of biology itself. Batch-correction isotropy therefore removes genuine biology in direct proportion to the non-exchangeability of batches, and the known-rare-to-unknown-rare transfer that grounds our validation (Lemma~\ref{lem:exchangeability}) is least guaranteed precisely where it is most clinically consequential. This is not a peripheral concern: the populations whose presence is diagnostic of response to immunotherapy or targeted therapy --- TPEX, LAMP3$^{+}$ mregDC, TREM2$^{+}$ LAM, CCR8$^{+}$ eTreg, the p-EMT hybrid, apCAF --- are all born at this junction. A framework that claims pan-cancer-atlas scalability must therefore not only satisfy the minimax envelope of \Cref{thm:minimax} but explicitly disclose the cancer-specific boundary condition under which its exchangeability bound is least tight.
